\section{Relações e funções}\label{sec:relacoes}
Nesta seção, introduziremos formalmente os conceitos básicos relacionados a relações e funções.
Oportunamente, introduziremos o Axioma da Potência.
\subsection{Relações}
\begin{definition}
    Uma \textbf{relação} é um conjunto de pares ordenados.
    Formalmente,

    \[
        \forall R\, \bigl(R \text{ é uma relação} \leftrightarrow \forall z \in R\, z \text{ é um par ordenado}\bigr).
    \]
\end{definition}

Conceitos fundamentais associados a relações são os de domínio e imagem.
Enunciaremos suas definições formais, e, a seguir, mostraremos que a existência de tais conjuntos pode ser justificada com os axiomas já introduzidos.

\begin{definition}
    Seja $R$ um conjunto.
    O domínio de $R$, denotado por $\dom R$, é o conjunto $z$ cujos elementos são as primeiras coordenadas dos pares ordenados de $R$.
    Formalmente,
    \[
        \forall z\,\bigl(z = \dom R \leftrightarrow \forall t\,(t \in z \leftrightarrow \exists u\, (t,u) \in R)\bigr).
    \]

    A imagem de $R$, denotada por $\ran R$, é o conjunto $z$ cujos elementos são as segundas coordenadas dos pares ordenados de $R$.
    Formalmente,
    \[
        \forall z\,\bigl(z = \ran R \leftrightarrow \forall t\,(t \in z \leftrightarrow \exists u\,(u,t) \in R)\bigr).
    \]
\end{definition}
\begin{lemma}
    Na notação acima, os conjuntos $\dom R$ e $\ran R$ estão bem definidos.
    Ou seja,
    \[
        \forall R\,\bigl(\exists! z\, \forall t\,(t \in z \leftrightarrow \exists u\,(t,u) \in R)\bigr).
    \]

    \[
        \forall R\,\bigl(\exists! z\, \forall t\,(t \in z \leftrightarrow \exists u\,(u,t) \in R)\bigr).
    \]
\end{lemma}
\begin{proof}
    Seja $R$ um conjunto. Se $u \in R$ é um par ordenado, então existem $a, b$ tais que $u = (a, b)=\{\{a\}, \{a, b\}\}$.
    Assim,
    \[
        \{a, b\} \in u \in R \implies \{a, b\} \in \bigcup R \implies \bigl(a \in \bigcup\bigcup R\bigr) \land \bigl(b \in \bigcup\bigcup R\bigr).
    \]

    Seja $z=\{a \in \bigcup\bigcup R : \exists u\,(a, u) \in R\}$.
    Vamos verificar que $z$ é tal que $\forall t\,(t \in z \leftrightarrow \exists u\,(t, u) \in R)$.

    Ora, dado $t \in z$, temos que $t \in \bigcup\bigcup R$ e $\exists u\,(t, u) \in R$ pela definição de $z$.
    Reciprocamente, se $\exists u\,(t, u) \in R$, então, pelo exposto acima, $t \in \bigcup\bigcup R$, e, portanto, $t \in z$.

    Da mesma forma, sendo $z'=\{b \in \bigcup\bigcup R : \exists u\,(u, b) \in R\}$, temos que $\forall t\,(t \in z' \leftrightarrow \exists u\,(u, t) \in R)$.

    Note que a existência de $z$ e $z'$ decorre do Esquema da Compreensão aplicado ao conjunto $\bigcup\bigcup R$.
    A unicidade decorre do Axioma da Extensão.
\end{proof}

\begin{definition}
    Dado um conjunto $A$, dizemos que uma relação $R$ é uma \emph{relação em $A$} se $\dom R \subseteq A$ e $\ran R \subseteq A$.
\end{definition}
\subsection{O Axioma da Potência e o produto Cartesiano}
Dados conjuntos $A$ e $B$, é natural considerar o conjunto de todos os pares ordenados $(a,b)$ tais que $a \in A$ e $b \in B$.
Quando existe, esse conjunto é usualmente chamado de produto cartesiano de $A$ e $B$.
É possível demonstrar que não é possível provar que tal relação existe para todos os conjuntos $A$ e $B$ utilizando apenas os axiomas já introduzidos.
Para garantir tal existência, introduziremos o Axioma da Potência.
Uma alternativa seria utilizar o Esquema da Substituição, mas optaremos por adiar a introdução desse axioma para um capítulo posterior.
\begin{axiom}[Axioma da Potência]
    Para qualquer conjunto $x$, existe um conjunto $y$ tal que todo subconjunto de $x$ é elemento de $y$.
    Formalmente,
    \[
        \forall x\, \exists y\, \forall z\,(z \subseteq x \rightarrow z \in y).
    \]
\end{axiom}

Utilizando uma instância do esquema da compreensão, dado um conjunto $x$, pode-se restringir o conjunto $y$ fornecido pelo Axioma da Potência para obter o conjunto cujos elementos são exatamente os subconjuntos de $x$.
Deixamos os detalhes ao leitor (Exercício~\ref{exercicio:existenciaPotencia}).

\begin{definition}
    Dado um conjunto $x$, o conjunto $\mathcal P(x)$ é o conjunto cujos elementos são exatamente os subconjuntos de $x$.
    Formalmente,
\[
    \forall y\,\bigl(y = \mathcal P(x) \leftrightarrow \forall z\,(z \in y \leftrightarrow z \subseteq x)\bigr).
\]
\end{definition}

Utilizando o conjunto das partes, podemos garantir a existência do produto cartesiano de quaisquer conjuntos $A$ e $B$, conforme exposto a seguir.
\begin{lemma}
    Dados conjuntos $A$ e $B$, existe um único conjunto $C$ que tem como elementos exatamente os pares ordenados $(a, b)$ tais que $a \in A$ e $b \in B$.
    Formalmente,
    \[
        \forall A\, \forall B\, \exists! C\, \forall z\,\bigl(z \in C \leftrightarrow \exists a\in A\, \exists b\in B\,(z = (a, b))\bigr).
    \]
\end{lemma}
\begin{proof}
    Dados $a \in A$ e $b \in B$, temos que $\{a\}, \{a, b\} \in \mathcal P(A\cup B)$, e, portanto, $(a, b) = \{\{a\}, \{a, b\}\} \in \mathcal P(\mathcal P(A\cup B))$.
    Seja $C = \{z \in \mathcal P(\mathcal P(A\cup B)) : \exists a\in A\, \exists b\in B\,(z = (a, b))\}$.

    Verifiquemos que $C$ é tal que $\forall z\,\bigl(z \in C \leftrightarrow \exists a\in A\, \exists b\in B\,(z = (a, b))\bigr)$.
    Dado $z \in C$, temos que $z \in \mathcal P(\mathcal P(A\cup B))$ e $\exists a\in A\, \exists b\in B\,(z = (a, b))$ pela definição de $C$.
    Reciprocamente, se $\exists a\in A\, \exists b\in B\,(z = (a, b))$, então, pelo exposto acima, $z \in \mathcal P(\mathcal P(A\cup B))$, e, portanto, $z \in C$.

    A existência de $C$ decorre do Esquema da Compreensão aplicado ao conjunto $\mathcal P(\mathcal P(A\cup B))$.
    A unicidade decorre do Axioma da Extensão.
\end{proof}

\begin{definition}
    Dado quaisquer conjuntos $A$ e $B$, o produto cartesiano de $A$ e $B$, denotado por $A \times B$, é o conjunto cujos elementos são exatamente os pares ordenados $(a, b)$ tais que $a \in A$ e $b \in B$.
    Formalmente,
    \[
        \forall C\,\bigl(C = A \times B \leftrightarrow \forall z\,(z \in C \leftrightarrow \exists a\in A\, \exists b\in B\,(z = (a, b)))\bigr).
    \]
\end{definition}

Outra operação natural associada a uma relação é a formação de sua inversa.
Intuitivamente, a relação inversa de $R$ é obtida trocando-se a ordem das coordenadas em cada par ordenado pertencente a $R$.

Por simplicidade, definiremos $R^{-1}$ para qualquer conjunto, independentemente de $R$ ser ou não uma relação.
Porém, raramente tal notação será utilizada no caso de $R$ não ser uma relação.

\begin{definition}
    Seja $R$ um conjunto.
    Definimos $R^{-1}$ como o conjunto dos pares ordenados $(b,a)$ tais que $(a,b) \in R$.
    Formalmente,
    \[
        \forall z\,\bigl(z = R^{-1} \leftrightarrow \forall t\,(t \in z \leftrightarrow \exists a\,\exists b\,((a,b)\in R \land t=(b,a)))\bigr).
    \]
\end{definition}

A escrita explícita acima é demasiadamente enfadonha e pouco natural na prática matemática.
A notação abaixo é mais natural, porém ainda não foi introduzida formalmente:
\[
    \{(b,a)\in \ran R\times \dom R : (a,b)\in R\}.
\]

Vamos aproveitar a oportunidade para introduzi-la.

\begin{definition}
    Considere $\varphi(a,b,A,B,t_1,\dots,t_n)$ uma fórmula.
    Dados conjuntos $A$ e $B$ e parâmetros $t_1, \dots, t_n$, definimos
    \[
        \{(a,b)\in A\times B : \varphi(a,b,A,B,t_1,\dots,t_n)\}
    \]
    como o conjunto
    \[
        \{c \in A\times B : \exists u\,\exists v\,(c=(u,v) \land \varphi(u,v,A,B,t_1,\dots,t_n))\}.
    \]
\end{definition}

\begin{lemma}
    Na notação acima, o conjunto $R^{-1}$ está bem definido.
    Ou seja,
    \[
        \forall R\,\exists! z\,\forall t\,(t \in z \leftrightarrow \exists a\,\exists b\,((a,b)\in R \land t=(b,a))).
    \]
\end{lemma}
\begin{proof}
    Seja $R$ um conjunto.
    Considere
    \[
        z=\{(u,v)\in \ran R\times \dom R : (v,u)\in R\}.
    \]

    Vamos verificar que
    \[
        \forall t\,(t \in z \leftrightarrow \exists a\,\exists b\,((a,b)\in R \land t=(b,a))).
    \]

    Suponha que $t \in z$.
    Pela definição da notação introduzida acima, existem $u$ e $v$ tais que
    \[
        t=(u,v) \land (v,u)\in R.
    \]
    Tomando $a=v$ e $b=u$, obtemos
    \[
        (a,b)\in R \quad\text{e}\quad t=(b,a).
    \]

    Reciprocamente, suponha que existam $a$ e $b$ tais que
    \[
        (a,b)\in R \quad\text{e}\quad t=(b,a).
    \]
    Então $a \in \dom R$ e $b \in \ran R$.
    Logo,
    \[
        (b,a)\in \ran R\times \dom R.
    \]
    Além disso, como $(a,b)\in R$, segue que
    \[
        t=(b,a)\in \{(u,v)\in \ran R\times \dom R : (v,u)\in R\}=z.
    \]

    A unicidade decorre do Axioma da Extensão.
\end{proof}

Podemos também definir a composição de relações.

\begin{definition}
    Sejam $R$ e $S$ relações.
    A composição de $R$ e $S$, denotada por $R\circ S$, é definida por
    \[
        R\circ S = \{(a, c)\in \dom S \times \ran R : \exists b\,((a, b) \in S \land (b, c) \in R)\}.
    \]
\end{definition}
\subsection{Funções}
Conforme usualmente definido, uma função é um tipo especial de relação.

\begin{definition}
    Uma \textbf{função} é uma relação $f$ tal que, para cada $a \in \dom f$, existe um único $b \in \ran f$ tal que $(a, b) \in f$.
    Formalmente,
    \[
        \forall f\,\bigl(f \text{ é uma função} \leftrightarrow f \text{ é uma relação} \land \forall a\in\dom f\,\exists! b\in\ran f\,((a,b)\in f)\bigr).
    \]
\end{definition}

Conforme usualmente feito, definiremos a notação $f(a)$ para denotar o único $b$ tal que $(a, b) \in f$, quando $f$ é uma função e $a \in \dom f$.

\begin{definition}
    Sejam $f$ uma função e $a \in \dom f$.
    O valor de $f$ em $a$, denotado por $f(a)$, é o único $b$ tal que $(a,b)\in f$.

    Em símbolos,
    \[
        \forall b\,\bigl(f(a) = b \leftrightarrow (a, b) \in f\bigr).
    \]
\end{definition}

Convém formalizar a notação usual $f:A\rightarrow B$.

\begin{definition}
    Sejam $f$ uma função e $A$ e $B$ conjuntos.
    \begin{enumerate}[label=(\alph*)]
        \item Dizemos que $f$ é uma função de $A$ em $B$, o que denotamos por $f:A\rightarrow B$, se $\dom f = A$ e $\ran f \subseteq B$.
        \item Dizemos que $f$ é uma função de $A$ se $\dom f = A$.
        \item Dizemos que $f$ é uma função parcial de $A$ se $\dom f \subseteq A$.
        \item Dizemos que $f$ é uma função em $B$ se $\ran f \subseteq B$.
        \item Dizemos que $f$ é uma função sobre $B$ se $\ran f = B$.
        \item Dizemos que $B$ é um contradomínio para $f$ se $\ran f \subseteq B$.
    \end{enumerate}
\end{definition}

As notações de imagem inversa, direta e de restrição também são as usuais.

\begin{definition}
    Sejam $f$ uma função, $A$ um conjunto e $B$ um conjunto.
    \begin{enumerate}[label=(\alph*)]
        \item A imagem direta de $A$ por $f$, denotada por $f[A]$, é o conjunto $\{b \in \ran f : \exists a \in A\, (f(a) = b)\}$.
        \item A imagem inversa de $B$ por $f$, denotada por $f^{-1}[B]$, é o conjunto $\{a \in \dom f : f(a) \in B\}$.
        \item A restrição de $f$ a $A$, denotada por $f\restriction A$, é a função cujos elementos são exatamente os pares $(a,b)\in f$ tais que $a\in A$.
    \end{enumerate}
\end{definition}

Notemos que as notações $f[A]$ e $f^{-1}[B]$ e $f\restriction A$ estão bem definidas mesmo que $A$ não seja um subconjunto de $\dom f$ ou $B$ não seja um subconjunto de $\ran f$.
Isso não é um problema, e, na verdade, é desejável em muitos momentos.
    
\subsection*{Exercícios}
\begin{exercise}\label{exercicio:existenciaPotencia}
    Justifique a existência de $\mathcal P(x)$.
    Ou seja, mostre que
    \[
        \forall x\, \exists! y\, \forall z\,(z \in y \leftrightarrow z \subseteq x).
    \]
\end{exercise}
\begin{exercise}
    Sejam $R$ e $S$ relações.
    Prove as seguintes afirmações.
    \begin{enumerate}[label=(\alph*)]
        \item $\dom R^{-1} = \ran R$;
        \item $\ran R^{-1} = \dom R$;
        \item $(R^{-1})^{-1} = R$;
        \item $\dom(R\circ S) \subseteq \dom S$;
        \item $\ran(R\circ S) \subseteq \ran R$;
        \item $(R\circ S)^{-1} = S^{-1}\circ R^{-1}$.
    \end{enumerate}
\end{exercise}

\begin{exercise}
    Sejam $f:A\rightarrow B$ e $C$ um conjunto.
    Mostre que $\dom (f\restriction C) = A\cap C$.
    Em particular, mostre que se $C \subseteq A$, então $\dom (f\restriction C) = C$.
\end{exercise}